\section{Measurement Protocol}
\label{sec:protocol}

\paragraph{Prompting.} Each vignette is presented to each model as an
independent single-turn call. The system prompt frames the task as
weighing in on a relationship dispute between ``Agent 1'' and ``Agent 2''
and asks the model to judge Agent 1's degree of fault, returning a single
JSON object with four fields: a short free-text \texttt{reasoning}, a
short \texttt{obligation\_identified} phrase naming what the model
understood Agent 1's responsibility to have been (a manipulation check --
did the model identify the intended norm, independent of how it rated
fault), an integer \texttt{fault\_rating}, and a self-reported
\texttt{confidence} field retained in the prompt for completeness but not
used as the confidence measure (see below). The user message is the
vignette text verbatim, with nothing added or wrapped around it. The
prompt itself never mentions gender, asks the model to consider it, or
asks the model to comment on it; the only place gender appears anywhere
in the input is the vignette's own explicit \texttt{(female)}/\texttt{(male)}
tag (Section~\ref{sec:dataset}). This is deliberate: a well-aligned model
asked directly whether gender affected its judgment will simply say no
regardless of what it actually did, so the only valid way to detect an
effect is by comparing separately elicited gender-configuration cells of
the same scenario/severity cell \emph{after} generation, never through
self-report on the sensitive attribute itself.

\paragraph{Fault rating scale.} \texttt{fault\_rating} is a single integer
0--7 scale with defined anchors at every point (0: ``not at all in the
wrong;'' 7: ``entirely in the wrong; a severe breach that would justify
significant relationship consequences''), rather than a separate
categorical verdict paired with a blameworthiness score. A categorical
verdict can still be derived post hoc via a cutpoint decided during
analysis (e.g.\ $\le 3$ as not-in-the-wrong-leaning, $\ge 4$ as
in-the-wrong-leaning) without asking the model for a category directly.
Each provider's native structured-output or tool-use feature is used to
enforce the output schema where available, rather than relying on
prompt-only compliance, since prompt-only JSON formatting is unreliable
across models from multiple providers; a defined fallback parser and
malformed-output rule apply where native schema enforcement is
unavailable. The decomposition into \texttt{obligation\_identified} (what
the model understood the obligation to be) and \texttt{fault\_rating}
(how much that obligation's violation counts against the agent) mirrors
the causal/responsibility/blame distinction established in the marital
attribution literature \citep{bradburyfincham1990} -- the closest
psychological-construct precedent for decomposing partner-directed fault
judgment this way, rather than treating blame as a single undifferentiated
intensity.

\paragraph{Confidence as dispersion, not self-report.} Confidence is not
treated as a prompt field to trust. It is measured empirically in two
passes: a \emph{confirmatory pass} (each of the 288 core vignettes, run
once per model at low/near-zero temperature, the primary data for RQ1--RQ2)
and a \emph{stability pass} (each vignette run $N$ additional times at a
higher, deployment-realistic temperature). The primary confidence measure
is the dispersion of \texttt{fault\_rating} across the $N$ stability-pass
samples (e.g.\ standard deviation) rather than a ``verdict flip rate,''
which was the metric in an earlier draft of this protocol but does not
transfer cleanly to a continuous rating -- a rating that moves from 3 to 4
across samples is not a ``flip'' the way a discrete category is. A
post-hoc categorical flip rate can still be reported as a secondary,
more interpretable measure derived from the same cutpoint used for the
categorical verdict above. Temperature-sensitivity findings in the
LLM-as-judge reliability literature motivate treating $N$ and temperature
as decisions to justify rather than default assumptions: same-verdict
consistency across repeated samples has been reported to fall
substantially as temperature rises from near-zero, and the temperature
that maximizes reliability is not universal across models or tasks
\citep{schroeder2024trust}, while a large-scale audit across judges and
benchmarks finds internal consistency itself is an unreliable proxy for
validity \citep{norman2026reliability} -- a caution directly relevant to
treating dispersion as a \emph{confidence} measure rather than mistaking
low dispersion for a guarantee that the rating is measuring the intended
construct. \textbf{Open calibration decisions, not yet finalized:} the
exact value of $N$ and the stability-pass sampling temperature, and
whether the stability pass runs on all 288 core vignettes or a
representative subset (running it on all 288 multiplies API calls by $N$
and is a real cost, not a rounding error).

\paragraph{Hedges, refusals, and malformed output.} Some fraction of
responses are expected to hedge (e.g.\ ``both are somewhat at fault''),
decline to commit to a rating, or fail schema validation outright. Per
RQ3, this is tracked as a variable of direct interest rather than
discarded as noise: whether a model hedges more often for a given agent
gender or norm family is treated as a legitimate finding about
differential willingness to render judgment, not a nuisance to code away.
The exact coding rule for hedges (e.g.\ whether the free-text
\texttt{reasoning} field is parsed for an implied lean, or counted as
missing data) is calibrated from a small pilot run on the drafted
scenarios before the confirmatory pass, rather than fixed in advance of
seeing how often it actually occurs.

\paragraph{Model roster.} Five models across four providers are evaluated:
Anthropic Claude Sonnet 5, OpenAI GPT-5-mini, Google Gemini 2.5 Flash,
Meta Llama 3.3 70B, and DeepSeek V3.2, accessed via a single inference
API (OpenRouter) for a consistent calling interface across providers.
